\uuid{KccK}
\exo7id{1850}
\titre{exo7 1850}
\auteur{ridde}
\organisation{exo7}
\datecreate{1999-11-01}
\isIndication{false}
\isCorrection{false}
\chapitre{Equation aux dérivées partielles}
\sousChapitre{Equation aux dérivées partielles}
\cpge{mpsi-mp2i.fonctions-deux-variables.b, pcsi.fonctions-deux-variables.b, ptsi.fonctions-deux-variables.b, mpsi-mp2i.fonctions-deux-variables.c, pcsi.fonctions-deux-variables.c, ptsi.fonctions-deux-variables.c, bcpst1.derivees-primitives-integrales.b, bcpst1.fonctions-deux-variables.c, mp-mpi.calcul-differentiel-optimisation.a}
\module{Equation différentielle}
\niveau{L3}
\difficulte{}

\contenu{
\question{
Soit $f : \Rr^3 \rightarrow \Rr$ d\'erivable. On pose $F (x, y, z) = f (x-y, y-z, z-x)$.\\
Calculer $\frac{\partial F}{\partial x} + \frac{\partial F}{\partial y} +
\frac{\partial F}{\partial z}$.
}
}
